\chapter{\texorpdfstring{\code{pdvar}}{pdvar}: Decision Matrix Functions}
\label{sec:pdvar}

\code{pdvar} creates continuous cell-local Bernstein decision expressions backed by YALMIP \code{sdpvar} coefficients. It participates in supported affine and known-data algebra. It does not choose solvers or call \code{optimize}. Comparison overloads create \code{pdlmi} objects, whose \code{toYalmip} method later converts them to YALMIP constraints. Shared matrix-shape, unary, structural, reduction, reordering, and repetition methods are implemented centrally by \code{pdbase} and remain complete public \code{pdvar} behavior; this chapter documents their decision-expression syntax, stable shape examples, validations, and rate-aware boundaries.

A degree-\(m\) \code{pdvar} represents a decision such as \(P(\boldsymbol\rho)\) or \(y_k(\boldsymbol\rho)\). For \(m\ge1\), neighboring cells reuse every decision handle on their complete shared face, which enforces value continuity without auxiliary equality constraints; for \(m=0\), the entire grid reuses one decision matrix. \code{rhodiff} applies cellwise forward differences scaled by \(1/h_s^{\mathbf c}\) and stores one row per vertex in \(\mathcal V_{\mathcal R}\), while \code{value} recovers numeric coefficient payloads only after the caller has validated YALMIP's solver status. Products remain affine only when the other factor is known numeric or \code{pdmat} data.

\section{API Map}
\label{sec:pdvar-lookup}

\begin{tabularx}{\textwidth}{@{}lY@{}}
\toprule
Task & Entry \\
\midrule
Construct decisions &
\hyperref[sec:pdvar-construction]{\code{pdvar}} \\
Inspect coefficients &
\hyperref[sec:pdvar-table]{\code{bernsteinTable}},
\hyperref[sec:pdvar-backend-inheritance]{shared storage/evaluation/elevation} \\
Convert assigned coefficients &
\hyperref[sec:pdvar-value]{\code{value}} \\
Differentiate rates &
\hyperref[sec:pdvar-rhodiff]{\code{rhodiff}} \\
Affine/mixed algebra &
\hyperref[sec:pdvar-algebra]{\code{+}},
\hyperref[sec:pdvar-algebra]{\code{-}},
\hyperref[sec:pdvar-algebra]{\code{*}} with known data \\
Matrix/index methods &
\hyperref[sec:pdvar-matrix-ops]{structural methods},
\hyperref[sec:pdvar-indexing]{indexing and assignment} \\
Build LMIs &
\hyperref[sec:pdvar-comparison]{\code{<=}},
\hyperref[sec:pdvar-comparison]{\code{>=}},
\hyperref[sec:pdvar-comparison]{\code{==}} \\
\bottomrule
\end{tabularx}

\section{Construct \texorpdfstring{\code{pdvar}}{pdvar} Decisions}
\index{affine decision expression}\index{decision matrix}
\label{sec:pdvar-construction}

\syntaxhead
\begin{lstlisting}[style=dpsyntax]
P = pdvar(n, gridVectors)
P = pdvar(n, n, gridVectors)
P = pdvar(n, p, gridVectors)
P = pdvar(..., "full")
P = pdvar(..., "symmetric")
P = pdvar(..., Degree=0)
P = pdvar(..., Degree=1)
P = pdvar(..., Degree=degree)
P = pdvar(..., RateBounds=rb)
P = pdvar(..., ValidationMode=mode)
\end{lstlisting}

\arghead
\begin{itemize}[leftmargin=*]
  \item \code{pdvar(n, gridVectors)} creates an \(n\times n\) square variable using scalar-size shorthand; for example, \code{pdvar(2, {[0 1]})}. Square variables default to symmetric YALMIP coefficients. \item \code{pdvar(n, n, gridVectors)} is the explicit square-size form; for example, \code{pdvar(2, 2, {[0 1]})}. It has the same default symmetric structure as the shorthand. \item \code{pdvar(n, p, gridVectors)} creates an \(n\times p\) variable. Rectangular variables default to full YALMIP coefficients; for example, \code{pdvar(2, 3, {[0 1]})}. \item \code{"full"} and \code{"symmetric"} explicitly choose the YALMIP coefficient structure. \code{"symmetric"} requires a square size. \item \code{Degree=degree} accepts one finite nonnegative integer scalar shorthand or an \(\ell\)-element vector and is stored as a \(1\)-by-\(\ell\) row \(\mathbf m=(m_1,\ldots,m_\ell)\). The omitted default is \(\mathbf1\). An explicitly supplied scalar on a multidimensional grid expands uniformly and emits \codeerr{pdvar}{ScalarDegreeExpansion} once; omitted defaults and one-dimensional forms remain warning-free. An all-zero vector creates one parameter-independent decision matrix shared across the grid. Otherwise neighboring cells share complete control faces; a zero component means constancy in that parameter direction. \item \code{RateBounds=rb} stores parameter-rate metadata for later \code{rhodiff(P)} calls. Use \code{RateBounds=[-1 1]} for one parameter; use a two-row matrix for two parameters, as shown in the rate-metadata example below. \item \code{ValidationMode=mode} accepts scalar text \code{"fast"} or \code{"strict"}, case-insensitively, and defaults to \code{"fast"}. It affects repeated validation of package-generated storage only; public options, grid, structure, affine safety, continuity, and source semantics remain global in both modes. The mode is not stored.
\end{itemize}

\begin{center}
\begin{tikzpicture}[node distance=8mm and 9mm,>=Latex,
  c/.style={dp/result,
    font=\scriptsize,align=center,minimum width=20mm,minimum height=8mm}]
  \node[c] (p10) {\(P^{(1)}[0]\)};
  \node[c,right=of p10] (p11) {\(P^{(1)}[1]\)};
  \node[c,right=of p11] (p12) {\(P^{(1)}[2]\)};
  \node[dp/label,above=of p11] {cell 1: degree-two local row};
  \draw[thick,dpblue] (p10) -- (p11) -- (p12);
  \node[c,below=of p12] (p20) {\(P^{(2)}[0]\)};
  \node[c,right=of p20] (p21) {\(P^{(2)}[1]\)};
  \node[c,right=of p21] (p22) {\(P^{(2)}[2]\)};
  \node[dp/label,above=of p21] {cell 2: degree-two local row};
  \draw[thick,dpblue] (p20) -- (p21) -- (p22);
  \draw[dp/arrow,<->] (p12.south) -- (p20.north);
  \node[dp/label,below=of p21,align=center]
    {same YALMIP handle: \(P^{(1)}[2]=P^{(2)}[0]\)};
\end{tikzpicture}

\textit{For positive degree, neighboring physical cells reuse complete boundary-face YALMIP handles; no extra continuity LMI is generated.}
\end{center}

\examplehead
The square shorthand creates a \(2\times2\) continuous decision object.  The default degree is one and the object carries no rate metadata.

\begin{lstlisting}[style=dpmatlab]
yalmip('clear')
Ps = pdvar(2, {[0 1]});
P= Ps.coeffs([1]);
symmetricCoefficient = ishermitian(P{1})
variableClass = class(Ps)
matrixSize = size(Ps)
degree = Ps.Degree
rateDependent = Ps.HasRateDependence
\end{lstlisting}

\begin{lstlisting}[style=dpoutput]
symmetricCoefficient = 1
variableClass = 'pdvar'
matrixSize = 1×2 double
     2     2

degree = 1

rateDependent =
  logical
   0
\end{lstlisting}

Use two dimension arguments for a rectangular full variable.

\begin{lstlisting}[style=dpmatlab]
yalmip('clear')
Pf = pdvar(2, 3, {[0 1]});
variableClass = class(Pf)
matrixSize = size(Pf)
degree = Pf.Degree
\end{lstlisting}

\begin{lstlisting}[style=dpoutput]
variableClass = 'pdvar'
matrixSize = 1×2 double
     2     3

degree = 1
\end{lstlisting}

The structure flag makes the coefficient structure explicit.  This is useful when a square variable should use full, not symmetric, coefficients.

\begin{lstlisting}[style=dpmatlab]
yalmip('clear')
Pfull = pdvar(2, 2, {[0 1]}, "full");
P = Pfull.coeffs([1]);
fullCoefficient = ishermitian(P{1})
variableClass = class(Pfull)
matrixSize = Pfull.MatrixSize
degree = Pfull.Degree
\end{lstlisting}

\begin{lstlisting}[style=dpoutput]
fullCoefficient = logical
   0

variableClass = 'pdvar'
matrixSize = 1×2 double
     2     2

degree = 1
\end{lstlisting}

\shortpairspace
Use \code{Degree=0} for a parameter-independent decision variable stored with grid metadata.


\begin{lstlisting}[style=dpmatlab]
yalmip('clear')
P0 = pdvar(1, [0 1 2], Degree=0);
variableClass = class(P0)
degree = P0.Degree
cellCount = P0.ncell()
\end{lstlisting}

\begin{lstlisting}[style=dpoutput]
variableClass = 'pdvar'
degree = 0
cellCount = 2
\end{lstlisting}

\shortpairspace
Direction-wise degree permits variation in selected parameters and exact constancy in the others.

\begin{lstlisting}[style=dpmatlab]
yalmip('clear')
Pa = pdvar(1, {[0 1], [10 20]}, Degree=[2 0]);
degree = Pa.Degree
labels = Pa.lbls()
coefficientCount = Pa.ncoeff()
\end{lstlisting}

\begin{lstlisting}[style=dpoutput]
degree = 1×2 double
     2     0

labels = 3×2 double
     0     0
     1     0
     2     0

coefficientCount = 3
\end{lstlisting}

\shortpairspace
Rate bounds are metadata for later \code{rhodiff(P)} calls.  \code{pdlmi} constrains rate vertices only after an expression such as \code{rhodiff} has stored rate-vertex rows.

\begin{lstlisting}[style=dpmatlab]
yalmip('clear')
Pr = pdvar(1, {[0 1], [10 20]}, RateBounds=[-1 2; -3 4]);
variableClass = class(Pr)
rateDependent = Pr.HasRateDependence
rateBounds = Pr.RateBounds
\end{lstlisting}

\begin{lstlisting}[style=dpoutput]
variableClass = 'pdvar'
rateDependent =
  logical
   1

rateBounds = 2×2 double
     -1     2
     -3     4
\end{lstlisting}

\remarkshead
\begin{itemize}[leftmargin=*]
  \item Missing or misplaced dimensions and grids raise \codeerr{pdvar}{InvalidInput}; invalid sizes and degrees raise \codeerr{pdvar}{Invalid}\allowbreak\code{MatrixSize} or \codeerr{pdvar}{InvalidDegree}. Explicit multidimensional scalar degree expansion is accepted with warning \codeerr{pdvar}{ScalarDegreeExpansion}. \item Option failures use \codeerr{pdvar}{Invalid}\allowbreak\code{Options}, \codeerr{pdvar}{Unknown}\allowbreak\code{Option}, or \codeerr{pdvar}{Unsupported}\allowbreak\code{Option}; invalid modes raise \codeerr{pdvar}{Invalid}\allowbreak\code{ValidationMode}. \item A symmetric nonsquare request raises \codeerr{pdvar}{InvalidStructure}; grid failures use \codeerr{pdvar}{InvalidGrid} or \codeerr{pdvar}{Invalid}\allowbreak\code{GridVector}. \item \code{RateBounds} must be a finite \(\ell\)-by-2 matrix with each lower bound no greater than its upper bound; violations raise \codeerr{pdbase}{Invalid}\allowbreak\code{RateBounds}.
\end{itemize}

\section{Inspect Object Properties}

\code{pdvar} declares no additional class-owned properties.  It exposes the read-only inherited \code{pdbase} properties listed in \cref{sec:pdbase-fields}; their private set access prevents user assignment. In particular, decision coefficients are inspected through \code{P.LocalValues}, while \code{P.Degree}, \code{P.IsContinuous}, and \code{P.RateBounds} describe their representation and rate metadata.

\begin{lstlisting}[style=dpmatlab]
yalmip('clear')
P = pdvar(2, {[0 1]}, "symmetric", Degree=2, RateBounds=[-1 1]);
degree = P.Degree
isContinuous = P.IsContinuous
rateBounds = P.RateBounds
firstCellCoefficientCount = numel(P.coeffs(1))
\end{lstlisting}

\begin{lstlisting}[style=dpoutput]
degree = 2

isContinuous =
  logical
   1

rateBounds = 1×2 double
     -1     1

firstCellCoefficientCount = 3
\end{lstlisting}

Use public algebra to create a modified value; direct assignments such as \code{P.Degree = 1} are not supported.

\section{Inspect, Evaluate, And Elevate Inherited Storage}
\label{sec:pdvar-backend-inheritance}

\syntaxhead
\begin{lstlisting}[style=dpsyntax]
cellSubs = cells(P)
labels = lbls(P)
coefficients = coeffs(P, cellSubs)
n = ncell(P)
n = ncoeff(P)
n = npar(P)
X = evaluate(P, point)
X = P.evaluate(point)
vals = P.elevVals(degreeIncrement)
Q = P.elevate(degreeIncrement)
\end{lstlisting}

\descriptionhead
The storage queries use the shared \code{pdbase} cell and label order and return the symbolic \code{sdpvar} coefficient payloads of the selected decision expression. \code{evaluate} reconstructs the affine symbolic matrix at one in-bounds parameter point; a derivative/rate-row object returns a \(1\)-by-\(2^\ell\) cell array in stored rate-vertex order. \code{elevVals} returns only an elevated nested coefficient tree. \code{elevate} returns a new \code{pdvar} of the same dynamic class with the same represented expression, physical cells, YALMIP decision variables, metadata, and rate-row order at the higher degree; it does not create new decision freedom or modify \code{P}. The backend performs the same binomially normalized \(\ell\)-dimensional convolution as known-data elevation, using the fixed separable kernel \(K_{\mathbf M-\mathbf m}[\mathbf r]=\prod_s\binom{M_s-m_s}{r_s}\). \Cref{sec:bernstein-degree-elevation} derives this exact representation change.

\examplehead
\begin{lstlisting}[style=dpmatlab]
yalmip('clear')
P = pdvar(2, {[0 1]}, "full", Degree=1);
physicalCells = P.cells()
labels = P.lbls()
coefficientTableSize = size(P.coeffs(1))
counts = [P.ncell(), P.ncoeff(), P.npar()]
X = P.evaluate(0.25);
Q = P.elevate(2);
sampleClass = class(X)
sampleSize = size(X)
elevatedClass = class(Q)
degrees = [P.Degree, Q.Degree]
elevatedCoefficientTableSize = size(Q.coeffs(1))
\end{lstlisting}
\begin{lstlisting}[style=dpoutput]
physicalCells = 1
labels = 2×1 double
     0
     1
coefficientTableSize = 1×2 double
     1     2
counts = 1×3 double
     1     2     1
sampleClass = 'sdpvar'
sampleSize = 1×2 double
     2     2
elevatedClass = 'pdvar'
degrees = 1×2 double
     1     3
elevatedCoefficientTableSize = 1×2 double
     1     4
\end{lstlisting}

Vector increments may elevate different axes by different amounts, including an axis on which the original decision is constant.

\begin{lstlisting}[style=dpmatlab]
yalmip('clear')
P = pdvar(1, {[0 1], [0 1]}, Degree=[1 0]);
Q = P.elevate([1 2]);
sourceDegree = P.Degree
elevatedDegree = Q.Degree
coefficientCounts = [P.ncoeff(), Q.ncoeff()]
elevatedCoefficientTableSize = size(Q.coeffs([1 1]))
\end{lstlisting}

\begin{lstlisting}[style=dpoutput]
sourceDegree = 1×2 double
     1     0

elevatedDegree = 1×2 double
     2     2

coefficientCounts = 1×2 double
     2     9

elevatedCoefficientTableSize = 1×2 double
     1     9
\end{lstlisting}

\remarkshead
Invalid physical-cell selectors raise \codeerr{pdbase}{InvalidCellSubs}. Malformed points raise \codeerr{pdvar}{InvalidPoint}; out-of-bounds points raise \codeerr{pdvar}{PointOutOfBounds}. Invalid elevation increments raise \codeerr{pdbase}{InvalidDegreeIncrement}. Evaluation returns symbolic expressions until the user assigns or solves the underlying YALMIP variables; use \code{value} for the post-assignment numeric conversion documented in \cref{sec:pdvar-value}.

\section{\texorpdfstring{\code{rhodiff}}{rhodiff}}
\index{rhodiff@\texttt{rhodiff}}
\label{sec:pdvar-rhodiff}

\syntaxhead
\begin{lstlisting}[style=dpsyntax]
D = rhodiff(P)
D = rhodiff(P, rb)
\end{lstlisting}

\arghead
\begin{itemize}[leftmargin=*]
  \item \code{P} is a \code{pdvar} without existing rate-vertex rows. \item \code{rb} is a finite \(\ell\)-by-\(2\) numeric matrix of lower and upper parameter-rate bounds, with each lower bound no greater than its upper bound. It must equal \code{P.RateBounds} when the object already carries nonempty bounds. \item Without \code{rb}, \code{rhodiff(P)} requires nonempty \code{P.RateBounds}. \item \code{D} is a discontinuous \code{pdvar} with one coefficient row per rate vertex. In one parameter, degree \(m>0\) becomes \(m-1\), while degree zero remains zero. In two or more parameters, every nonzero-axis partial is elevated exactly from \(\mathbf m-\mathbf e_s\) to the common tensor degree \(\mathbf m=\code{P.Degree}\) before summation; a zero-degree axis contributes zero. An all-zero tensor degree returns a complete zero row for every rate vertex.
\end{itemize}

For a coefficient row of direction-wise degree \(\mathbf m\) on physical cell \(\mathcal H_{\mathbf c}\), define \(h_s^{\mathbf c}=\rho_s^{(c_s+1)}-\rho_s^{(c_s)}>0\). Before rate substitution, differentiation along a nonzero-degree axis uses the forward difference
\[
\left(\frac{\partial P^{(\mathbf c)}}{\partial\rho_s}\right)[\mathbf i]
=\frac{m_s}{h_s^{\mathbf c}}
\left(P^{(\mathbf c)}[\mathbf i+\mathbf e_s]-P^{(\mathbf c)}[\mathbf i]\right),
\qquad i_s=0,\ldots,m_s-1.
\]
Thus the \(s\)-th partial has degree \(\mathbf m-\mathbf e_s\). In one parameter this reduced degree is the public result degree. In two or more parameters, \code{rhodiff} elevates every nonzero partial back to the common degree \(\mathbf m\) before adding the rate-weighted rows, so partials from different axes have compatible tensor labels. \Cref{sec:bernstein-differentiation} derives the reduction and exact re-elevation.

For a one-parameter degree-one coefficient pair \(P^{(c)}[0],P^{(c)}[1]\), the rate-dependent derivative row reduces to
\[
  \dot P(\rho,\nu)
=\nu\,\frac{P^{(c)}[1]-P^{(c)}[0]}{h_1^c}.
\]
This is the one-dimensional special case: the two rows correspond to the lower and upper allowed scheduling rates.

\begin{center}
\begin{tikzpicture}[>=Latex,node distance=7mm and 10mm,
  box/.style={dp/neutral,font=\scriptsize,
    align=center,minimum width=33mm,minimum height=10mm},
  result/.style={dp/result,
    font=\scriptsize,align=center,minimum width=34mm,minimum height=10mm}]
  \node[box] (p) {\(P^{(c)}[0],P^{(c)}[1]\)\\one-parameter cell \(c\)};
  \node[result,right=of p] (d)
    {\((P^{(c)}[1]-P^{(c)}[0])/h_1^c\)\\local derivative};
  \node[result,right=of d,yshift=8mm] (lo) {\(\underline\nu\) row};
  \node[result,right=of d,yshift=-8mm] (hi) {\(\overline\nu\) row};
  \draw[->,thick,dpblue] (p) -- (d);
  \draw[->,thick,dpblue] (d) -- (lo);
  \draw[->,thick,dpblue] (d) -- (hi);
\end{tikzpicture}

\textit{Differentiation changes local coefficients; rate substitution creates rows.  Neither operation adds physical cells.}
\end{center}

For \(\ell\) parameters, nonuniform grids use a distinct step length in every direction and cell:
\[
    \dot P^{(\mathbf c)}(\boldsymbol\rho,\boldsymbol\nu)=\sum_{s=1}^{\ell}\nu_s\,\frac{\partial P^{(\mathbf c)}}{\partial\rho_s}(\boldsymbol\rho).
\]
\code{RateBounds} fixes \(\mathcal R\), so \code{rhodiff} evaluates this affine expression at the \(2^\ell\) vertices in \(\mathcal V_{\mathcal R}\) and stores one row per vertex. Supply those bounds explicitly as \code{rhodiff(P, rb)}, or initialize \code{P} with \code{RateBounds=rb}; if both are present they must match. Without bounds there is no finite vertex set to store, so \code{rhodiff(P)} raises \codeerr{pdvar}{MissingRateBounds}. \Cref{sec:bernstein-rate-vertex-storage} gives the exact hypercube-local row and coefficient layout.

\remarkshead
Unsupported call arity or differentiating an expression that already contains rate-vertex rows raises \codeerr{pdvar}{InvalidDiff}.  Explicit malformed rate bounds raise \codeerr{pdvar}{Invalid}\allowbreak\code{RateBounds}.  Calling \code{rhodiff(P)} without stored bounds raises \codeerr{pdvar}{Missing}\allowbreak\code{RateBounds}; explicit bounds that disagree with \code{P.RateBounds} raise \codeerr{pdvar}{RateBounds}\allowbreak\code{Mismatch}.

\examplehead
When rate bounds live on the object, \code{rhodiff(P)} uses them directly.

\begin{lstlisting}[style=dpmatlab]
yalmip('clear')
P = pdvar(2, {[0 1 2]}, "symmetric", RateBounds=[-1 1]);
D = rhodiff(P);
derivativeClass = class(D)
derivativeDegree = D.Degree
isContinuous = D.IsContinuous
coefficientTableSize = size(D.coeffs(1))
rateBounds = D.RateBounds
\end{lstlisting}

\begin{lstlisting}[style=dpoutput]
derivativeClass = 'pdvar'
derivativeDegree = 0
isContinuous =
  logical
   0
coefficientTableSize = 1×2 double
     2     1
rateBounds = 1×2 double
     -1     1
\end{lstlisting}

For a two-parameter object, the rate-vertex rows enumerate all lower/upper combinations from the two rows of \code{RateBounds}.

\begin{lstlisting}[style=dpmatlab]
yalmip('clear')
Q = pdvar(1, {[0 1], [10 20]}, RateBounds=[-1 1; -2 2]);
E = rhodiff(Q);
% Rows 1:4 use [-1 -2], [-1 2], [1 -2], and [1 2], respectively.
derivativeClass = class(E)
derivativeDegree = E.Degree
coefficientTableSize = size(E.coeffs([1 1]))
rateVertexCount = coefficientTableSize(1)
coefficientsPerRateVertex = coefficientTableSize(2)
rateBounds = E.RateBounds
\end{lstlisting}

\begin{lstlisting}[style=dpoutput]
derivativeClass = 'pdvar'
derivativeDegree = 1
coefficientTableSize = 1×2 double
     4     4
rateVertexCount = 4
coefficientsPerRateVertex = 4
rateBounds = 2×2 double
     -1     1
     -2     2
\end{lstlisting}

A zero-degree axis contributes exactly zero but remains present in the common multivariate degree and in the rate-vertex enumeration.

\begin{lstlisting}[style=dpmatlab]
yalmip('clear')
P = pdvar(1, {[0 1], [10 20]}, Degree=[2 0]);
D = rhodiff(P, [-1 1; -2 2]);
derivativeDegree = D.Degree
coefficientTableSize = size(D.coeffs([1 1]))
rateVertexCount = coefficientTableSize(1)
coefficientsPerRateVertex = coefficientTableSize(2)
\end{lstlisting}

\begin{lstlisting}[style=dpoutput]
derivativeDegree = 1×2 double
     2     0

coefficientTableSize = 1×2 double
     4     3

rateVertexCount = 4
coefficientsPerRateVertex = 3
\end{lstlisting}

\section{\texorpdfstring{\code{bernsteinTable}}{bernsteinTable}}
\label{sec:pdvar-table}

\syntaxhead
\begin{lstlisting}[style=dpsyntax]
T = bernsteinTable(P)
T = bernsteinTable(P, cellSubs)
T = bernsteinTable(P, "oneLine")
T = bernsteinTable(P, cellSubs, "oneLine")
\end{lstlisting}

\arghead
\begin{itemize}[leftmargin=*]
  \item \code{P} is a \code{pdvar} object. \item \code{cellSubs} selects one physical cell; omit it to list every cell. \item \code{"oneLine"} is the only supported text option.  It returns one compact Bernstein expression per selected cell. \item \code{T} is a table.  Ordinary rows contain the local label, basis, physical-node flag, and symbolic or numeric value.  A rate-dependent derivative also includes \code{RateVertexIndex} and \code{RateVertex}.
\end{itemize}

\begin{center}
\begin{tikzpicture}[>=Latex,node distance=6mm and 9mm,
  box/.style={dp/neutral,font=\scriptsize,
    align=center,minimum width=31mm,minimum height=9mm},
  result/.style={dp/result,
    font=\scriptsize,align=center,minimum width=39mm,minimum height=11mm}]
  \node[box] (ordinary) {ordinary \code{pdvar}\\cell + local label};
  \node[box,below=of ordinary] (rate) {\code{rhodiff(P)}\\cell + rate vertex + label};
  \node[result,right=of ordinary,yshift=-7mm] (table) {stable display text\\
    never raw \code{sdpvar} cells};
  \draw[->,thick,dpblue] (ordinary.east) -- (table.west);
  \draw[->,thick,dpblue] (rate.east) -- (table.west);
\end{tikzpicture}
\end{center}

\examplehead
\begin{lstlisting}[style=dpmatlab]
yalmip('clear')
P = pdvar(1, {[0 1]});
T = bernsteinTable(P, "oneLine");
tableHeight = height(T)
tableWidth = width(T)
hasExpression = ismember("Expression", string(T.Properties.VariableNames))
\end{lstlisting}

\begin{lstlisting}[style=dpoutput]
tableHeight = 1
tableWidth = 2
hasExpression = logical
   1
\end{lstlisting}

\shortpairspace
For \code{D = rhodiff(P)}, \code{bernsteinTable(D, cellSubs, "oneLine")} emits one row for every selected physical cell and rate-box vertex, rather than one row per local Bernstein coefficient.  \code{RateVertexIndex} and \code{RateVertex} identify the lower/upper rate tuple; their tensor order is defined in \Cref{sec:bernstein-rate-vertex-storage}.

\begin{lstlisting}[style=dpmatlab]
yalmip('clear')
P = pdvar(1, {[0 1], [0 1]}, RateBounds=[-1 1; -2 2]);
D = rhodiff(P);
T = bernsteinTable(D, [1 1], "oneLine");
vertexIndices = T.RateVertexIndex
rateVertices = vertcat(T.RateVertex{:})
expressionCount = height(T)
\end{lstlisting}

\begin{lstlisting}[style=dpoutput]
vertexIndices = 4×1 double
     1
     2
     3
     4

rateVertices = 4×2 double
     -1    -2
     -1     2
     1    -2
     1     2

expressionCount = 4
\end{lstlisting}

\remarkshead
\code{bernsteinTable} is an inspection method, not a solver call.  Text options other than \code{"oneLine"}, repeated selectors, invalid physical-cell selectors, or inconsistent rate-vertex rows raise \codeerr{pdvar}{Invalid}\allowbreak\code{BernsteinTableInput}.

\section{Build Affine Algebra with Known Data}
\label{sec:pdvar-algebra}

A \code{pdvar} result must remain affine in the underlying YALMIP decisions. The following blocks separate affine sums from products so the supported boundary is visible beside each example.

\subsection{Signs, Addition, And Subtraction}

Use this family to combine compatible decision expressions, coefficient-backed known data, finite real numeric data, and affine \code{sdpvar} matrices. For two parameter-dependent operands, alignment proceeds in a fixed order: verify matching outer bounds, form the sorted union of interior nodes on each axis, exactly subdivide/re-express both coefficient trees on that common grid, elevate both to the componentwise maximum degree, and finally add or subtract matching coefficients. Numeric and affine constant operands are promoted before the last two steps. The corresponding overload names are \code{plus}, \code{minus}, \code{uplus}, and \code{uminus}; see \cref{sec:pdbase-mergegrid,sec:bernstein-degree-elevation}.

\syntaxhead
\begin{lstlisting}[style=dpsyntax]
R = P + Q
R = P + A
R = A + P
R = P + M
R = M + P
R = P + X
R = X + P
R = P - Q
R = P - A
R = A - P
R = P - M
R = M - P
R = +P
R = -P
\end{lstlisting}

Here \code{A} is coefficient-backed \code{pdmat}, \code{M} is numeric, and \code{X} is affine \code{sdpvar}.  A numeric scalar broadcasts to the payload size; an affine \code{sdpvar} is promoted as parameter-independent data.

\examplehead
\begin{lstlisting}[style=dpmatlab]
yalmip('clear')
P = pdvar(2, {[0 1]}, "full");
X = sdpvar(2, 2, 'full');
S = P + X;
D = eye(2) - P;
N = -P;
sumClass = class(S)
sumSize = size(S)
sumDegree = S.Degree
sumContainsDecision = S.ContainsDecision
differenceClass = class(D)
negativeClass = class(N)
\end{lstlisting}

\begin{lstlisting}[style=dpoutput]
sumClass = 'pdvar'
sumSize = 1×2 double
     2     2
sumDegree = 1
sumContainsDecision =
  logical
   1
differenceClass = 'pdvar'
negativeClass = 'pdvar'
\end{lstlisting}

A nonlinear symbolic operand such as \code{eta*eta} is rejected. Rate-vertex rows from \code{rhodiff} may be added to compatible ordinary expressions; the ordinary rows are broadcast over the active rate vertices.

The next example exercises both stages of parameter-dependent alignment: a coarse degree-one decision is refined to the known operand's two-cell grid and elevated to degree two before coefficientwise addition.

\begin{lstlisting}[style=dpmatlab]
yalmip('clear')
P = pdvar(1, {[0 1]}, Degree=1);
A = pdmat({[0 0.5 1]}, @(rho) rho.^2, Degree=2);
S = P + A;
resultClass = class(S)
mergedGrid = S.GridInfo.Vectors{1}
resultDegree = S.Degree
storageCounts = [S.ncell(), S.ncoeff()]
\end{lstlisting}

\begin{lstlisting}[style=dpoutput]
resultClass = 'pdvar'

mergedGrid = 1×3 double
         0    0.5000    1.0000

resultDegree = 2

storageCounts = 1×2 double
     2     3
\end{lstlisting}

\subsection{Multiplication By Known Or Numeric Data}
\label{sec:pdvar-operator-entries}

The \code{mtimes} overload supports multiplication only when decision dependence occurs on at most one side. A coefficient-backed \code{pdmat} or finite real numeric matrix may multiply a \code{pdvar} from either side. A scalar \code{pdvar} may scale a known matrix. Two decision-bearing factors, including \code{P*Q} and \code{P*eta} for a decision \code{sdpvar eta}, are nonlinear and rejected.

\syntaxhead
\begin{lstlisting}[style=dpsyntax]
R = A * P
R = P * A
R = M * P
R = P * M
\end{lstlisting}

The first example shows numeric left and right matrix products without exposing incidental YALMIP variable identifiers.

\examplehead
\begin{lstlisting}[style=dpmatlab]
yalmip('clear')
P = pdvar(2, 1, {[0 1]}, "full");
L = [1 2] * P;
R = P * [4 5];
S = 3 * P;
leftSize = size(L)
rightSize = size(R)
scaledSize = size(S)
degrees = [L.Degree R.Degree S.Degree]
\end{lstlisting}

\begin{lstlisting}[style=dpoutput]
leftSize = 1×2 double
     1     1
rightSize = 1×2 double
     2     2
scaledSize = 1×2 double
     2     1
degrees = 1×3 double
     1     1     1
\end{lstlisting}

A known parameter-dependent factor adds its Bernstein degree componentwise through ordered coefficient convolution. On one parameter cell, a quadratic decision and affine known factor produce
\[
(PA)[k]=\sum_{i+j=k}
  \frac{\binom2i\binom1j}{\binom3k}\,P[i]A[j],
  \qquad k=0,1,2,3.
\]
The same backend rule is documented at \code{bernProd} in \cref{sec:pdbase-bernprod} and derived step by step in \cref{sec:bernstein-product-math}. In \(\ell\) parameters, the ordered products are grouped by \(\mathbf k=\mathbf i+\mathbf j\) and weighted by the product of the one-dimensional binomial normalization ratios; the result degree is the componentwise sum.

\begin{lstlisting}[style=dpmatlab]
yalmip('clear')
P = pdvar(2, {[0 1]}, "symmetric", Degree=2);
A = pdmat({[0 1]}, {eye(2), 2*eye(2)}, Degree=1);
L = A * P;
R = P * A;
knownLeftClass = class(L)
knownLeftDegree = L.Degree
knownRightClass = class(R)
knownRightDegree = R.Degree
coefficientsPerCell = numel(L.coeffs(1))
\end{lstlisting}

\begin{lstlisting}[style=dpoutput]
knownLeftClass = 'pdvar'
knownLeftDegree = 3
knownRightClass = 'pdvar'
knownRightDegree = 3
coefficientsPerCell = 4
\end{lstlisting}

Direction-wise degrees add independently. A known factor that varies only in the second parameter raises only the second component of the result degree.

\begin{lstlisting}[style=dpmatlab]
yalmip('clear')
P = pdvar(1, {[0 1], [0 1]}, Degree=[2 0]);
A = pdmat({[0 1], [0 1]}, ...
    @(rho1, rho2) 1 + rho2, Degree=[0 1]);
L = A * P;
R = P * A;
leftDegree = L.Degree
rightDegree = R.Degree
coefficientsPerCell = L.ncoeff()
\end{lstlisting}

\begin{lstlisting}[style=dpoutput]
leftDegree = 1×2 double
     2     1

rightDegree = 1×2 double
     2     1

coefficientsPerCell = 6
\end{lstlisting}

A derivative object may be multiplied by numeric data or a matching-grid known object while preserving every rate row.  Products with rate dependence on both sides, products of a derivative with an ordinary decision expression, and products involving metadata-only rate dependence are rejected.

Two decision-bearing factors are also rejected: their product would be quadratic in YALMIP decision coefficients.  The supported boundary is numeric or known \code{pdmat} data multiplying one affine \code{pdvar} expression on either side, with the written matrix order preserved.

\remarkshead
Incompatible affine operands raise \codeerr{pdvar}{InvalidAddition} or \codeerr{pdvar}{InvalidSubtraction}.  Bad dimensions, two decision-bearing factors, nonlinear symbolic data, or unsupported rate combinations raise \codeerr{pdvar}{InvalidMultiplication}.  Incompatible parameter bounds raise \codeerr{pdvar}{MixedGrid}; a function-only \code{pdmat} operand raises \codeerr{pdvar}{FunctionOnlyAlgebra}.

\section{Reshape and Transform Decision Matrices}
\label{sec:pdvar-matrix-ops}

Every method below is inherited from the centralized \code{pdbase} implementation and transforms each local affine YALMIP payload while the \code{pdvar} reconstruction hook preserves the physical grid, degree, class, decision metadata, and compatible rate rows. The examples report stable classes, dimensions, and metadata rather than YALMIP's internal variable numbers.

\subsection{Shape, Equality, And MATLAB Protocols}
\label{sec:pdvar-matrix-shape}
\label{sec:pdvar-matrix-protocol}

Shape queries describe the matrix payload.  \code{isequal} is an exact identity check: grid, degree, matrix size, metadata, and local symbolic payloads must already match; it does not merge or elevate operands.

\syntaxhead
\begin{lstlisting}[style=dpsyntax]
sz = size(P)
d = size(P, dim)
n = height(P)
n = width(P)
n = length(P)
n = numel(P)
n = ndims(P)
tf = isequal(P, Q, ...)
idx = end(P, k, n)
n = numArgumentsFromSubscript(P, S, ctx)
\end{lstlisting}

\examplehead
\begin{lstlisting}[style=dpmatlab]
yalmip('clear')
P = pdvar(2, 3, {[0 1]}, "full");
matrixSize = size(P)
rowCount = height(P)
columnCount = width(P)
longestDimension = length(P)
elementCount = numel(P)
dimensionCount = ndims(P)
equalsUnaryPlus = isequal(P,+P)
\end{lstlisting}

\begin{lstlisting}[style=dpoutput]
matrixSize = 1×2 double
     2     3

rowCount = 2
columnCount = 3
longestDimension = 3
elementCount = 6
dimensionCount = 2
equalsUnaryPlus =
  logical
   1
\end{lstlisting}

\code{end}, \code{subsref}, \code{subsasgn}, and \code{numArgumentsFromSubscript} support the two-index and dot-access forms shown in \cref{sec:pdvar-indexing}; they do not expose physical-cell storage.

\subsection{Transpose And Conjugate Transpose}
\label{sec:pdvar-matrix-transform}

Both forms transpose every local affine payload without breaking the YALMIP coefficient graph.

\syntaxhead
\begin{lstlisting}[style=dpsyntax]
T = transpose(P)
H = ctranspose(P)
T = P.'
H = P'
\end{lstlisting}

\examplehead
\begin{lstlisting}[style=dpmatlab]
yalmip('clear')
P = pdvar(2, 3, {[0 1]}, "full");
T = P.';
H = P';
transposeSize = size(T)
conjugateTransposeSize = size(H)
transposeClass = class(T)
conjugateTransposeClass = class(H)
\end{lstlisting}

\begin{lstlisting}[style=dpoutput]
transposeSize = 1×2 double
     3     2
conjugateTransposeSize = 1×2 double
     3     2
transposeClass = 'pdvar'
conjugateTransposeClass = 'pdvar'
\end{lstlisting}

\begin{center}
\begin{tikzpicture}[>=Latex,node distance=28mm,
  a/.style={matrix of math nodes,nodes={draw=black!45,minimum size=6mm,
    font=\scriptsize},column sep=-\pgflinewidth,row sep=-\pgflinewidth}]
  \matrix[a] (m) {p_{11}&p_{12}&p_{13}\\p_{21}&p_{22}&p_{23}\\};
  \matrix[a,right=of m] (t) {p_{11}&p_{21}\\p_{12}&p_{22}\\p_{13}&p_{23}\\};
  \draw[->,thick,dpblue] (m.east) -- node[above,font=\scriptsize]
    {each coefficient} (t.west);
\end{tikzpicture}
\end{center}

\subsection{Vectorization, Reshape, And Squeeze}

\code{vec} uses MATLAB column-major order.  \code{reshape} accepts two positive dimensions with at most one inferred \code{[]} dimension and preserves the element count.  \code{squeeze} retains the two-dimensional package convention.

\syntaxhead
\begin{lstlisting}[style=dpsyntax]
V = vec(P)
R = reshape(P, [m n])
R = reshape(P, m, n)
S = squeeze(P)
\end{lstlisting}

\examplehead
\begin{lstlisting}[style=dpmatlab]
yalmip('clear')
P = pdvar(2, 3, {[0 1]}, "full");
V = vec(P);
R = reshape(P, [3 2]);
S = squeeze(P);
vectorSize = size(V)
reshapeSize = size(R)
squeezeSize = size(S)
\end{lstlisting}

\begin{lstlisting}[style=dpoutput]
vectorSize = 1×2 double
     6     1
reshapeSize = 1×2 double
     3     2
squeezeSize = 1×2 double
     2     3
\end{lstlisting}

Malformed or size-changing reshape requests raise \codeerr{pdvar}{InvalidReshape}.

\subsection{Diagonals And Trace}

\code{diag} extracts a diagonal or constructs a diagonal matrix from a vector; the selected offset may not be empty.  \code{trace} returns a scalar decision expression.

\syntaxhead
\begin{lstlisting}[style=dpsyntax]
d = diag(P)
d = diag(P, k)
D = diag(v)
D = diag(v, k)
t = trace(P)
\end{lstlisting}

\examplehead
\begin{lstlisting}[style=dpmatlab]
yalmip('clear')
P = pdvar(3, {[0 1]}, "full");
d = diag(P);
D = diag(d);
t = trace(P);
diagonalSize = size(d)
rebuiltSize = size(D)
traceSize = size(t)
\end{lstlisting}

\begin{lstlisting}[style=dpoutput]
diagonalSize = 1×2 double
     3     1
rebuiltSize = 1×2 double
     3     3
traceSize = 1×2 double
     1     1
\end{lstlisting}

Invalid offsets or empty selections raise \codeerr{pdvar}{InvalidDiag}.

\subsection{Triangular Parts And Reductions}
\label{sec:pdvar-matrix-reductions}

Triangular methods zero the complementary part of every coefficient. Reductions act on matrix dimensions; \code{"all"} is available for \code{sum} and \code{mean}.

\syntaxhead
\begin{lstlisting}[style=dpsyntax]
L = tril(P)
L = tril(P, k)
U = triu(P)
U = triu(P, k)
S = sum(P)
S = sum(P, dim)
S = sum(P, vecdim)
S = sum(P, "all")
M = mean(P)
M = mean(P, dim)
M = mean(P, vecdim)
M = mean(P, "all")
C = cumsum(P)
C = cumsum(P, dim)
C = cumsum(P, direction)
C = cumsum(P, dim, direction)
\end{lstlisting}

First form the two triangular views.

\examplehead
\begin{lstlisting}[style=dpmatlab]
yalmip('clear')
P = pdvar(3, {[0 1]}, "full");
L = tril(P);
U = triu(P);
lowerSize = size(L)
upperSize = size(U)
degrees = [L.Degree U.Degree]
\end{lstlisting}

\begin{lstlisting}[style=dpoutput]
lowerSize = 1×2 double
     3     3
upperSize = 1×2 double
     3     3
degrees = 1×2 double
     1     1
\end{lstlisting}

Then reduce a rectangular expression.

\begin{lstlisting}[style=dpmatlab]
yalmip('clear')
P = pdvar(2, 3, {[0 1]}, "full");
rowSum = sum(P, 2);
allMean = mean(P, "all");
running = cumsum(P, 2);
rowSumSize = size(rowSum)
allMeanSize = size(allMean)
runningSize = size(running)
\end{lstlisting}

\begin{lstlisting}[style=dpoutput]
rowSumSize = 1×2 double
     2     1
allMeanSize = 1×2 double
     1     1
runningSize = 1×2 double
     2     3
\end{lstlisting}

\code{vecdim} is a nonempty vector of unique positive integer dimensions; dimensions above two are singleton no-ops. \code{direction} is \code{"forward"} or \code{"reverse"}, case-insensitively. A direction supplied without a dimension uses the first nonsingleton matrix dimension, and \code{cumsum} along a dimension above two returns the same symbolic payloads. Invalid triangular offsets raise \codeerr{pdvar}{InvalidTriangularPart}. Invalid reduction calls raise \codeerr{pdvar}{InvalidSum}, \codeerr{pdvar}{InvalidMean}, or \codeerr{pdvar}{InvalidCumsum}; missing-value and native-output flags are not supported.

\subsection{Reordering And Rotation}
\label{sec:pdvar-matrix-reorder}

These methods rearrange matrix entries inside every symbolic coefficient; they do not reverse the parameter grid.

\syntaxhead
\begin{lstlisting}[style=dpsyntax]
Q = flip(P)
Q = flip(P, dim)
Q = flipud(P)
Q = fliplr(P)
Q = rot90(P)
Q = rot90(P, k)
\end{lstlisting}

\examplehead
\begin{lstlisting}[style=dpmatlab]
yalmip('clear')
P = pdvar(2, 3, {[0 1]}, "full");
F = flipud(P);
R = rot90(P);
flippedSize = size(F)
rotatedSize = size(R)
flippedClass = class(F)
rotatedClass = class(R)
\end{lstlisting}

\begin{lstlisting}[style=dpoutput]
flippedSize = 1×2 double
     2     3
rotatedSize = 1×2 double
     3     2
flippedClass = 'pdvar'
rotatedClass = 'pdvar'
\end{lstlisting}

\begin{center}
\begin{tikzpicture}[>=Latex,node distance=8mm,
  b/.style={dp/result,
    font=\scriptsize,align=center,minimum width=30mm,minimum height=12mm}]
  \node[b] (p) {\(2\times3\) decision matrix};
  \node[b,right=of p] (f) {\code{flipud}: \(2\times3\)};
  \node[b,right=of f] (r) {\code{rot90}: \(3\times2\)};
  \draw[->,thick,dpblue] (p) -- (f);
  \draw[->,thick,dpblue] (f) -- (r);
\end{tikzpicture}
\end{center}

Invalid calls raise \codeerr{pdvar}{InvalidFlip} or \codeerr{pdvar}{InvalidRot90}.

\subsection{Concatenation And Block Assembly}
\label{sec:pdvar-matrix-assembly}

Assembly accepts compatible \code{pdvar}, coefficient-backed \code{pdmat}, affine \code{sdpvar}, and finite real numeric blocks. It aligns grids and elevates all blocks exactly to the componentwise maximum degree while preserving an affine expression. \code{cat} supports only matrix dimensions 1 and 2.

\syntaxhead
\begin{lstlisting}[style=dpsyntax]
H = horzcat(P, Q, ...)
V = vertcat(P, Q, ...)
H = [P, Q]
V = [P; Q]
C = cat(dim, P, Q, ...)
D = blkdiag(P, Q, ...)
\end{lstlisting}

\examplehead
\begin{lstlisting}[style=dpmatlab]
yalmip('clear')
P = pdvar(2, {[0 1]}, "full");
H = [P, P];
V = [P; P];
D = blkdiag(P, eye(2));
horizontalSize = size(H)
verticalSize = size(V)
blockDiagonalSize = size(D)
\end{lstlisting}

\begin{lstlisting}[style=dpoutput]
horizontalSize = 1×2 double
     2     4
verticalSize = 1×2 double
     4     2
blockDiagonalSize = 1×2 double
     4     4
\end{lstlisting}

Incompatible blocks or rate metadata raise \codeerr{pdvar}{InvalidConcatenation}; unsupported dimensions raise \linebreak[4]\code{pdvar:}\allowbreak\code{UnsupportedCat}\allowbreak\code{Dimension}; invalid block-diagonal inputs raise \codeerr{pdvar}{InvalidBlkdiag}.  Incompatible bounds raise \codeerr{pdvar}{MixedGrid}; function-only known blocks raise \linebreak[4]\code{pdvar:}\allowbreak\code{FunctionOnly}\allowbreak\code{Algebra}.

\subsection{Repetition}

\code{repmat} tiles the matrix payload while retaining the parameter metadata. It accepts a positive scalar, two positive counts, or a two-element count vector.

\syntaxhead
\begin{lstlisting}[style=dpsyntax]
Q = repmat(P, m)
Q = repmat(P, m, n)
Q = repmat(P, [m n])
\end{lstlisting}

\examplehead
\begin{lstlisting}[style=dpmatlab]
yalmip('clear')
P = pdvar(1, 2, {[0 1]}, "full");
Q = repmat(P, [2 2]);
repeatedSize = size(Q)
repeatedDegree = Q.Degree
containsDecision = Q.ContainsDecision
\end{lstlisting}

\begin{lstlisting}[style=dpoutput]
repeatedSize = 1×2 double
     2     4
repeatedDegree = 1
containsDecision =
  logical
   1
\end{lstlisting}

Invalid repetition shapes raise \codeerr{pdvar}{InvalidRepmat}.

\section{Index and Assign \texorpdfstring{\code{pdvar}}{pdvar} Entries}
\label{sec:pdvar-indexing}

Parenthesis indexing selects a matrix block from every symbolic coefficient. Assignment writes a compatible numeric, \code{pdvar}, coefficient-backed \code{pdmat}, or affine \code{sdpvar} block into that position after the same common-grid refinement and componentwise maximum-degree elevation used by affine addition. The left object and parameter-dependent right-hand side are exactly re-expressed on the sorted union grid before the selected payload block is replaced at every aligned label. It does not select one physical cell. See \cref{sec:pdbase-mergegrid,sec:bernstein-degree-elevation}.

\syntaxhead
\begin{lstlisting}[style=dpsyntax]
Q = P(rows, cols)
Q = P(1:end, end)
value = P.PropertyName
P(rows, cols) = numericValue
P(rows, cols) = Q
P(rows, cols) = A
P(rows, cols) = X
\end{lstlisting}

\examplehead
\begin{lstlisting}[style=dpmatlab]
yalmip('clear')
P = pdvar(2, 3, {[0 1]}, "full");
lastCol = P(:, end);
P(:, 1) = 0;
selectedSize = size(lastCol)
assignedSize = size(P)
assignedClass = class(P)
assignedDegree = P.Degree
containsDecision = P.ContainsDecision
\end{lstlisting}

\begin{lstlisting}[style=dpoutput]
selectedSize = 1×2 double
     2     1
assignedSize = 1×2 double
     2     3
assignedClass = 'pdvar'
assignedDegree = 1
containsDecision =
  logical
   1
\end{lstlisting}

Assignment may therefore increase both the grid resolution and the degree of the complete decision expression while retaining its affine decision content.

\begin{lstlisting}[style=dpmatlab]
yalmip('clear')
P = pdvar(2, 2, {[0 1]}, "full", Degree=1);
A = pdmat({[0 0.5 1]}, ...
    @(rho) [rho.^2; 1+rho], Degree=2);
P(:, 2) = A;
assignedGrid = P.GridInfo.Vectors{1}
assignedDegree = P.Degree
assignedClass = class(P)
storageCounts = [P.ncell(), P.ncoeff()]
containsDecision = P.ContainsDecision
\end{lstlisting}

\begin{lstlisting}[style=dpoutput]
assignedGrid = 1×3 double
         0    0.5000    1.0000

assignedDegree = 2
assignedClass = 'pdvar'

storageCounts = 1×2 double
     2     3

containsDecision =
  logical
   1
\end{lstlisting}

Invalid selectors raise \codeerr{pdvar}{InvalidSubscript}.  Linear indexing, deletion, or non-parenthesis assignment raises \codeerr{pdvar}{UnsupportedAssignment}; incompatible or nonlinear right-hand sides raise \codeerr{pdvar}{InvalidAssignment}.
\section{\texorpdfstring{\code{value}}{value}}
\index{value@\texttt{value}}
\label{sec:pdvar-value}

\syntaxhead
\begin{lstlisting}[style=dpsyntax]
A = value(P)
rows = value(rhodiff(P))
\end{lstlisting}

\code{value} converts assigned \code{pdvar} coefficients into known, coefficient-backed \code{pdmat} data. An ordinary expression returns one \code{pdmat}, preserving its grid, matrix size, degree, local coefficient order, and continuity metadata. A rate-vertex expression created by \code{rhodiff} returns a \(1\)-by-\(2^\ell\) cell array of \code{pdmat} objects in the deterministic lower/upper vertex order induced by applying \code{helper.combRows} to \code{num2cell(RateBounds,2).'}. The returned \code{pdmat} objects do not carry rate metadata. Every symbolic coefficient must have an assigned finite real YALMIP value; otherwise the method raises \codeerr{pdvar}{UnassignedValue}.

\examplehead
\begin{lstlisting}[style=dpmatlab]
yalmip('clear')
P = pdvar(1, [0 1], Degree=2);
c = P.coeffs(1);
assign([c{:}], [1 2 4]);
A = value(P);
assignedCoefficients = cell2mat(A.coeffs(1))
\end{lstlisting}

\begin{lstlisting}[style=dpoutput]
assignedCoefficients = 1×3 double
     1     2     4
\end{lstlisting}

For rate rows, \code{value} returns one known object per rate vertex rather than collapsing the rows.

\begin{lstlisting}[style=dpmatlab]
yalmip('clear')
P = pdvar(1, [0 1], Degree=1);
c = P.coeffs(1);
assign([c{:}], [1 3]);
D = rhodiff(P, [-1 1]);
rows = value(D);
rowCount = numel(rows)
lowerRowClass = class(rows{1})
upperRowClass = class(rows{2})
lowerRateCoefficients = cell2mat(rows{1}.coeffs(1))
upperRateCoefficients = cell2mat(rows{2}.coeffs(1))
\end{lstlisting}

\begin{lstlisting}[style=dpoutput]
rowCount = 2
lowerRowClass = 'pdmat'
upperRowClass = 'pdmat'
lowerRateCoefficients = -2
upperRateCoefficients = 2
\end{lstlisting}

\section{Comparisons To \texorpdfstring{\code{pdlmi}}{pdlmi}}
\label{sec:pdvar-comparison}

The \code{le}, \code{ge}, and \code{eq} overloads implement \code{<=}, \code{>=}, and \code{==}.  Inequalities select either semidefinite or entry-wise meaning once for the complete original residual; equality is always direct coefficient equality.

\syntaxhead
\begin{lstlisting}[style=dpsyntax]
C = P <= 0
C = P >= 0
C = lhs <= rhs
C = lhs >= rhs
C = lhs == rhs
\end{lstlisting}

\arghead
\begin{itemize}[leftmargin=*]
  \item \code{rhs} may be compatible numeric, coefficient-backed \code{pdmat}, affine \code{sdpvar}, or \code{pdvar} data. Ordinary operands use the same grid refinement, componentwise maximum-degree alignment, and promotion as subtraction. \item For \code{<=} and \code{>=}, every coefficient in every physical cell and rate row is scanned before assembly. The complete relation is semidefinite only if every original coefficient is square Hermitian. Numeric Hermitian testing uses the inclusive rule \(\lVert A-A^{\mathsf T}\rVert_\infty\le10^{-10}\); symbolic coefficients use YALMIP's \code{ishermitian}. If any coefficient fails, the complete inequality is entry-wise and warning \codeerr{pdlmi}{ElementwiseInequality} is issued once. Modes are never mixed coefficient by coefficient or entry by entry. \item \code{==} accepts rectangular and nonsymmetric residuals without that warning. It equates corresponding coefficient entries directly. Ordinary rows compare with ordinary rows; derivative rows compare only with compatible derivative rows. \item \code{C} is a \code{pdlmi} wrapper. A comparison assembles constraints but does not solve an optimization problem.
\end{itemize}

\examplehead
Comparison creates a \code{pdlmi} wrapper with one stored constraint for each physical cell and local Bernstein coefficient.

\begin{lstlisting}[style=dpmatlab]
yalmip('clear')
P = pdvar(2, {[0 0.5 1]}, "symmetric");
Cneg = P <= 0;
negativeWrapperClass = class(Cneg)
negativeConstraintCount = numel(Cneg.Constraints)
Cpos = P >= 0;
positiveWrapperClass = class(Cpos)
positiveConstraintCount = numel(Cpos.Constraints)
\end{lstlisting}

\begin{lstlisting}[style=dpoutput]
negativeWrapperClass = 'pdlmi'
negativeConstraintCount = 4
positiveWrapperClass = 'pdlmi'
positiveConstraintCount = 4
\end{lstlisting}

A rectangular residual selects entry-wise inequality globally.  Equality is entry-wise by definition and is warning-free.

\begin{lstlisting}[style=dpmatlab]
yalmip('clear')
V = pdvar(3, 2, {[0 1]}, "full");
warnId = "pdlmi:ElementwiseInequality";
warnState = warning("off", warnId);
restoreWarning = onCleanup(@() warning(warnState));
Centry = V >= 0;
Ceq = V == sdpvar(3, 2, 'full');
entrywiseConstraintGroups = numel(Centry.Constraints)
equalityConstraintGroups = numel(Ceq.Constraints)
identity = V == V;
identityConstraints = identity.toYalmip()
clear restoreWarning
\end{lstlisting}

\begin{lstlisting}[style=dpoutput]
entrywiseConstraintGroups = 2
equalityConstraintGroups = 2

identityConstraints =

     []
\end{lstlisting}

Each group above corresponds to one local coefficient and contains six scalar entries in MATLAB column-major order. Under P\'olya, Putinar, SparseFullBox, or FullBox, each entry of an entry-wise inequality receives an independent scalar certificate. Equality is direct-only: constructor certificate options and \code{applyPolya}, \code{applyPutinar}, \code{applySparseFullBoxPreorder}, or \code{applyFullBoxPreorder} raise \codeerr{pdlmi}{Unsupported}\allowbreak\code{EqualityCertificate} before validating an increment, width, or order. Mixed ordinary/derivative row kinds raise \codeerr{pdvar}{InvalidEqualityRows}. Use overloaded \code{==} to build a modeling equality; use \code{isequal} only for MATLAB structural comparison.

\begin{center}
\begin{tikzpicture}[node distance=7mm and 12mm]
  \node[dp/decision,align=center,font=\scriptsize,text width=34mm,
    minimum height=11mm] (lhs) {\code{lhs}\\affine \code{pdvar}};
  \node[dp/known,align=center,font=\scriptsize,text width=34mm,
    minimum height=11mm,below=of lhs] (rhs)
    {\code{rhs}\\one or more compatible terms};
  \node[dp/contribution,right=of lhs,yshift=-7mm,align=center,
    font=\scriptsize,text width=30mm,minimum height=15mm] (res)
    {\code{lhs-rhs}\\stored residual};
  \node[dp/certificate,right=of res,align=center,font=\scriptsize,
    text width=27mm,minimum height=15mm] (relation)
    {relation\\\code{<=}, \code{>=}, or \code{==}};
  \node[dp/result,right=of relation,align=center,font=\scriptsize,
    text width=25mm,minimum height=15mm] (c)
    {\code{pdlmi}\\direct by default};
  \draw[dp/arrow] (lhs.east) -- ++(5mm,0) |- (res.west);
  \draw[dp/arrow] (rhs.east) -- ++(5mm,0) |- (res.west);
  \draw[dp/arrow] (res) -- (relation);
  \draw[dp/arrow] (relation) -- (c);
\end{tikzpicture}
\end{center}

\section{Boundaries and Related APIs}
\label{sec:pdvar-limits}

\begin{enumerate}[leftmargin=*]
  \item Constructor-created \code{pdvar} objects support every finite nonnegative scalar or direction-wise integer degree. Each cell stores \(\prod_s(m_s+1)\) Bernstein control matrices; a zero-degree axis is constant in that direction. \item Products may contain decision variables on at most one side; nonlinear decision products such as \code{P * Q} are outside this slice. \item Function-only \code{pdmat} operands cannot enter \code{pdvar} algebra unless they were constructed with explicit Bernstein evidence. \item \code{rhodiff} of an existing rate-vertex derivative expression is not part of the current public workflow.
\end{enumerate}

\paragraph{Related APIs.}
\code{pdvar}, \code{rhodiff}, \code{bernsteinTable}, \code{mtimes}, \code{pdlmi}, \code{pdmat}.
